\subsection*{4.2 Iterating Hensel lifting}
Let $f=gh$, where $g$ is irreducible and coprime to $h$, and $f,g,h$ are monic in $x$ with $x$-degree $\geq1$.
To start the Hensel lifting procedure, we factor the univariate polynomial $f(x,0)=f\bmod y$ as $f(x,0)=g_0(x)h_0(x)$, where $g_0$ is a divisor of $g\bmod y$, and coprime to $h_0$, and $\deg_x(g_0)\geq1$. Recall that by the pre-processing in Section 3.5, we may assume that there is such a decomposition of $f(x,0)$.
By the Euclidian algorithm, there are polynomials $a_0(x),b_0(x)$ such that $a_0g_0+b_0h_0=1$. Hence, for $I_0=\langle y\rangle$, we have $a_0g_0+b_0h_0\equiv1\pmod{I_0}$ and initiate Hensel lifting with
\[f\equiv g_0h_0\pmod{I_0}.\]
We iteratively apply Hensel lifting to $g_0,h_0$ as described in the proof of Lemma 12. Each time, the ideal gets squared. Let $I_k=I_0^{2^k}$. That is, we get polynomials $g_k,h_k$ such that
\[f\equiv g_kh_k\pmod{I_k},\]
and $g_k,h_k$ is a lift of $g_{k-1},h_{k-1}$ w.r.t.$f$. The following lemma states that $g_k$ divides $g$ modulo $I_k$, for all $k\geq0$. In a sense, the $g_k$'s approximate $g$ modulo increasing powers of $y$.
\begin{ntlemma}[Author Lemma 13]
With the notation from above, for all $k\geq0$ and some polynomial $h'_k$,
\[g\equiv g_kh'_k\pmod{I_k}\quad\text{and}\quad h_k\equiv hh'_k\pmod{I_k}.\]
Moreover,$g_k,h'_k$ is a lift of $g_{k-1},h'_{k-1}$ w.r.t.$g$ and $\deg_x(h'_k)\leq\deg_x(h_k)=2^{O(k)}$, for $k\geq1$.
\end{ntlemma}
\begin{proof}
The proof is by induction on $k\geq0$. For the base case, we have that $g_0$ divides $g$ modulo $I_0$, as explained above. Thus, for some polynomial $h'_0$ that is coprime to $g_0$, we have
\[g\equiv g_0h'_0\pmod{I_0}.\]
Hence, we have $h_0\equiv h'_0h\pmod{I_0}$. Note that $h'_0$ can be just1.
For the inductive step, assume that
\begin{equation}g\equiv g_{k-1}h'_{k-1}\pmod{I_{k-1}}\quad\text{and}\quad h_{k-1}\equiv hh'_{k-1}\pmod{I_{k-1}}.\tag{5}\end{equation}
Let $g'_k,h''_k$ be a lift of $g_{k-1},h'_{k-1}$ w.r.t.$g$, so that in particular
\begin{equation}g'_kh''_k\equiv g\pmod{I_k}.\tag{6}\end{equation}
We claim that then $g'_k,hh''_k$ is a lift of $g_{k-1},hh'_{k-1}$, i.e., of $g_{k-1},h_{k-1}$ by(5), w.r.t.$f$.
\begin{ntclaim}[Author Claim 14]
$g'_k,hh''_k$ is a lift of $g_{k-1},h_{k-1}$ w.r.t.$f$.
\end{ntclaim}
\begin{proof}
We check the three conditions for a lift in Definition 11. For the product condition(i), we have by(6)
\[g'_khh''_k=(g'_kh''_k)h\equiv gh\pmod{I_k}.\]
For condition(ii), we have $g'_k\equiv g_{k-1}\pmod{I_{k-1}}$ by assumption and similarly
\[hh''_k\equiv hh'_{k-1}\equiv h_{k-1}\pmod{I_{k-1}}.\]
By the remark after Definition 11, the condition(iii) already follows now. This proves Claim 14.
\end{proof}
Recall that also $g_k,h_k$ is a lift of $g_{k-1},h_{k-1}$. Hence, by the uniqueness property of Hensel lifting, there is a $u\in I_{k-1}$ such that
\begin{equation}g'_k\equiv g_k(1+u)\pmod{I_k}\quad\text{and}\quad hh''_k\equiv h_k(1-u)\pmod{I_k}.\tag{7}\end{equation}
Now observe that we can move the factor $1+u$: we have that $g_k(1+u),hh''_k$ is a lift of $g_{k-1},h_{k-1}$, then also $g_k,hh''_k(1+u)$ is a lift of $g_{k-1},h_{k-1}$.
\begin{ntclaim}[Author Claim 15]
$g_k,hh''_k(1+u)$ is a lift of $g_{k-1},h_{k-1}$ w.r.t.$f$.
\end{ntclaim}
\begin{proof}
We check the conditions for a lift in Definition 11. The first two of them are trivial: moving the factor $1+u$ clearly does not change the product. Because $u\in I_{k-1}$ we still have the equality with the factors $g_{k-1}$ and $h_{k-1}$ modulo $I_{k-1}$, respectively.
By the remark after Definition 11, the third condition already follows, but it also easy to check now:
Let $a,b\in R$ such that $ag_k+bh_k\equiv1\pmod{I_k}$. It follows by(7) that
\[ag_k+bhh''_k(1+u)\equiv ag_k+bh_k(1-u)(1+u)\equiv ag_k+bh_k(1-u^2)\equiv1\pmod{I_k}.\]
This proves Claim 15.
\end{proof}
Now, define $h'_k=h''_k(1+u)$. Note that
\begin{equation}h'_k\equiv h''_k\equiv h'_{k-1}\pmod{I_{k-1}}.\tag{8}\end{equation}
By(7), we have
\begin{equation}hh'_k\equiv hh''_k(1+u)\equiv h_k(1-u)(1+u)\equiv h_k\pmod{I_k}.\tag{9}\end{equation}
By(9) we have
\begin{equation}f=gh\equiv g_kh_k\equiv g_khh'_k\pmod{I_k}.\tag{10}\end{equation}
It follows from(10) that $gh\equiv g_kh'_kh\pmod{I_k}$. Now we want to cancel $h$ in the last equation and conclude that $g\equiv g_kh'_k\pmod{I_k}$. This we can do because $h$ is monic in $x$, it does not contain a factor $y$, i.e.$h\notin I_0$. Hence, together with(8), we conclude that $g_k,h'_k$ is a lift of $g_{k-1},h'_{k-1}$ w.r.t.$g$.
For the $x$-degree of $h'_k$, consider the equation $h_k\equiv hh'_k\pmod{I_k}$. Since $h$ is monic in $x$ the highest $x$-degree term in the product $hh'_k$ will survive the modulo operation. Therefore $\deg_x(h'_k)\leq\deg_x(h)+\deg_x(h'_k)=\deg_x(h_k)$.
To bound the degree of $h_k$ observe that in each round of the Hensel lifting, the maximum possible degree is bounded by a constant multiple($\leq5$) of the maximum degree in the previous round. After $k$ iterations, the degree of the lifted factors is therefore bounded by $2^{O(k)}$.
\end{proof}
For the ABP-size, we observed at the end of Section 4.1 that the size increases by a constant factor in one iteration. Hence, after $k$ iterations, the size increases by a factor $2^{O(k)}$.
